\subsection{Direct-kernel controls and component attribution}

The added controls use all six full molecular datasets and three seeds. A
direct-kernel oracle stores the float32 similarity matrix and supplies row-to-fold
mass, similarity columns, diagonal terms and cut evaluations to the unchanged
\texttt{balanced\_lloyd} and \texttt{fm\_polish} routines. The factor arm supplies
the same operations through compact coordinates. The two arms start from
identical seeded feasible assignments, permit at most 25 assignment iterations
and $n$ refinement operations, and use the same acceptance tests. We compare
ordinary and cubed Tanimoto similarities at exact counts and a 1\% corridor,
using 64, 256 and 512 landmarks. The dense arm is a heuristic reference, not an
optimal solution. Comparisons of a dense $BB^\top$ oracle with its factor verify
that the interface invokes equivalent optimization operations.

Separate 1- and 5-second controls include representation construction and use
the same ordered restart seeds. A supervising process stops the timed optimizer
at its deadline. Only checkpoints whose objective evaluation and label-array
serialization completed before the cutoff are eligible; the initial feasible
assignment is retained if construction exhausts the budget. Both arms repeatedly
call the current one-step assignment and refinement routines. Scheduling and
termination overshoot is recorded, and late work cannot improve the selected
partition. The supplement also retains an earlier cooperative protocol that
checks time between operations and allows the last operation to finish.

A separate 36-case grid measures dense-versus-factor cost on nested samples of
4,000, 16,000 and 32,000 ChEMBL fingerprints, using three seeds, ranks
64/256/512, ordinary similarity and a 1\% corridor. The fixed sampling order
matches the existing ChEMBL scaling protocol. Initialization, search limits,
two-core allocations, timing scope and independent full-kernel evaluation
match the iteration-controlled molecular study. These controls do not replace
the existing public-API scaling measurements with their different restart and
hardware settings.

Component controls retain the same feasible initialization when assignment or
refinement is omitted. Other rows vary one versus four restarts, uniform versus
similarity-spreading landmarks, and ordinary versus cubed kernels. The four-start
setting refines two candidates. Three additional rows share one capped group
array, a factor-independent feasible group packing, and strict preservation of
all groups. The grouped optimizer receives that packing as a warm start and one
additional seeded restart, and refines one candidate. Grouping changes the
feasible set, so its effect is distinguished from optimization within fixed
groups. Component timings disclose shared representation-construction costs.

Independent full NumPy float64 Tanimoto scoring evaluates every returned
partition. We report absolute cut error, error divided by total off-diagonal
kernel mass, and error divided by the true cut. Both denominators use the same
kernel power as the factor. Cut-relative error is undefined at a zero true cut;
it is not stabilized by adding an arbitrary constant. These representation
errors are distinct from the original-kernel objectives attained by the two
optimizers. The earlier 900 full-kernel comparisons are also reanalysed with
this cut-relative denominator. Independent RDKit bit-vector calculations
replay the added original-kernel scores and check saved counts, groups and
CCPart threshold edges.

\subsection{CCPart comparison and native constraints}

We use the unmodified CCPart implementation from Hestia-GOOD 1.2.2, upstream
commit \texttt{54b83e2}, with similarity thresholds 0.4, 0.6 and 0.8, requested
80/10/10 sizes and no response-label stratification. Every unordered Tanimoto
pair is considered when constructing the threshold graph; no top-$k$ truncation
is used. Native timing includes similarity construction, connected components
and packing. The three repetitions permute input order and invert that
permutation for evaluation. CCPart is deterministic: these repetitions measure
order sensitivity rather than stochastic optimizer variability.

We retain native counts, groups and all entities. Empty folds are reported
explicitly. When all folds are nonempty, two low-rank interventions use the
native assignment as a feasible initialization: one preserves only its exact
counts; the other also preserves every connected component. The latter uses
the same grouped warm-start budget as the component controls. Their timings
include factor construction and optimization but reuse the native groups;
they are optimization interventions rather than independently timed group
construction pipelines. The count-only intervention can split components;
the group-preserving intervention retains CCPart's threshold separation while
optimizing aggregate similarity. The full native size inventory accompanies
the quality and cost summaries.
